Fix an array in \(\Theta_{\mathrm{CTC}}(B)\) satisfying the three conditions in the statement. Membership in \(\mathrm{def{:}collision\text{-}model\text{-}class}\) supplies the assumptions cited below. Put
\[
U_{j,n}=U_{j,n}(t_n),\qquad
a_n=\beta_nU_{1,n},\qquad b_n=U_{2,n},\qquad M_n=\max(a_n,b_n).
\tag{1}
\]
By the definitions of \(c_{j,n}\),
\[
U_{j,n}^{\alpha_{j,n}}=c_{j,n}t_n.
\tag{2}
\]
Assumption \(\mathrm{ass{:}tail\text{-}constant\text{-}balance}\) gives \(\log(c_{1,n}/c_{2,n})=O(1)\), while \(\mathrm{ass{:}local\text{-}tail\text{-}collision}\) and the hypotheses give \(\Delta_n\to0\), \(\Delta_n\log t_n\to\lambda\), and \(\Delta_n\log c_{2,n}\to0\). Hence the exact identity
\[
\log\frac{U_{1,n}}{U_{2,n}}
=\frac{\log(c_{1,n}/c_{2,n})}{\alpha_{1,n}}
-\frac{\Delta_n(\log c_{2,n}+\log t_n)}{\alpha_{1,n}\alpha_{2,n}}
\tag{3}
\]
has the finite limit
\[
\log\frac{U_{1,n}}{U_{2,n}}
\longrightarrow \frac{\log c_0}{\alpha_0}-\frac{\lambda}{\alpha_0^2}.
\tag{4}
\]
Together with \(\beta_n\to\beta_0>0\) from \(\mathrm{ass{:}coefficient\text{-}limit}\), this proves that there is a fixed \(K<\infty\) such that
\[
K^{-1}\leq\frac{a_n}{b_n}\leq K
\tag{5}
\]
eventually. This is the required threshold-scale comparison.

Let \(q_{2,n}=F_{2,n}^{\leftarrow}(1-1/t_n)\). By \(\mathrm{ass{:}continuous\text{-}margins}\), \(P(X_{2,n}>q_{2,n})=1/t_n\), and conditioning on a marginal rank exceedance is almost surely equivalent to conditioning on the corresponding strict quantile exceedance. Assumption \(\mathrm{ass{:}local\text{-}tail\text{-}collision}\) also gives constants \(0<a_-<a_+<\infty\) such that
\[
a_-\leq\alpha_{j,n}\leq a_+\qquad(j=1,2)
\tag{6}
\]
eventually. Fix \(d\in(0,1)\). If \(M_n=a_n\), \(\mathrm{ass{:}uniform\text{-}tail\text{-}normalization}\) at the fixed multiple \(dU_{1,n}\) gives
\[
t_nP(\beta_n\epsilon_{1,n}>dM_n)=d^{-\alpha_{1,n}}{1+o(1)}>1
\tag{7}
\]
eventually; if \(M_n=b_n\), the same argument with the second innovation gives (7) with \(\alpha_{2,n}\). Since \(X_{2,n}=\beta_n\epsilon_{1,n}+\epsilon_{2,n}\) eventually by \(\mathrm{ass{:}linear\text{-}scm}\), positivity implies
\[
q_{2,n}\geq dM_n.
\tag{8}
\]
For the upper bound, choose a fixed \(C>2\) so large that \(2(C/2)^{-a_-}<1/2\). From (5), each of \(CM_n/(2a_n)\) and \(CM_n/(2b_n)\) lies in a fixed compact subset of \((0,\infty)\). Thus \(\mathrm{ass{:}uniform\text{-}tail\text{-}normalization}\), (6), and the union bound give
\[
\begin{split}
P(X_{2,n}>CM_n)
&\leq P(\beta_n\epsilon_{1,n}>CM_n/2)
+P(\epsilon_{2,n}>CM_n/2)\\
&<\frac1{t_n}
\end{split}
\tag{9}
\]
eventually. Therefore
\[
dM_n\leq q_{2,n}\leq CM_n.
\tag{10}
\]
In particular, \(q_{2,n}/a_n\) and \(q_{2,n}/b_n\) are bounded above and away from zero.

Define
\[
p_{1,n}=P(\beta_n\epsilon_{1,n}>q_{2,n}),
\qquad p_{2,n}=P(\epsilon_{2,n}>q_{2,n}).
\tag{11}
\]
Equations (2), (10), and \(\mathrm{ass{:}uniform\text{-}tail\text{-}normalization}\) imply
\[
p_{1,n}=\frac1{t_n}\left(\frac{q_{2,n}}{a_n}\right)^{-\alpha_{1,n}}\{1+o(1)\},
\qquad
p_{2,n}=\frac1{t_n}\left(\frac{q_{2,n}}{b_n}\right)^{-\alpha_{2,n}}\{1+o(1)\}.
\tag{12}
\]
Thus each \(p_{j,n}\) is of order \(1/t_n\), uniformly along this fixed array.

For completeness, we derive the needed one-large-jump decomposition rather than assume it. For fixed \(\delta\in(0,1/2)\), the same inclusions as in the exact-Pareto calculation give
\[
\begin{split}
p_{1,n}+p_{2,n}-p_{1,n}p_{2,n}
&\leq P(X_{2,n}>q_{2,n})\\
&\leq P(\beta_n\epsilon_{1,n}>(1-\delta)q_{2,n})
+P(\epsilon_{2,n}>(1-\delta)q_{2,n})\\
&\quad+P(\beta_n\epsilon_{1,n}>\delta q_{2,n})
P(\epsilon_{2,n}>\delta q_{2,n}),
\end{split}
\tag{13}
\]
where the product uses \(\mathrm{ass{:}innovation\text{-}independence}\). By (10), all arguments in (13), divided by their corresponding \(U_{j,n}\), remain in fixed compact subsets of \((0,\infty)\). Hence uniform tail normalization gives
\[
\frac{P(\beta_n\epsilon_{1,n}>(1-\delta)q_{2,n})}{p_{1,n}}
=(1-\delta)^{-\alpha_{1,n}}+o(1)
\tag{14}
\]
and the analogous second-component identity. The product term in (13) is \(O(t_n^{-2})=o(p_{1,n}+p_{2,n})\) by (12) and \(t_n\to\infty\), the latter following from \(\mathrm{ass{:}intermediate\text{-}sequence}\). Divide (13) by \(p_{1,n}+p_{2,n}\), first let \(n\to\infty\), and then let \(\delta\downarrow0\). Using (6) proves
\[
P(X_{2,n}>q_{2,n})=p_{1,n}+p_{2,n}+o(p_{1,n}+p_{2,n}).
\tag{15}
\]

Let \(H_n=F_{1,n}(X_{1,n})\). Forward orientation gives \(X_{1,n}=\epsilon_{1,n}\), and continuity gives \(E(H_n)=1/2\). On \(\{\beta_n\epsilon_{1,n}>q_{2,n}\}\),
\[
0\leq1-H_n\leq P(\epsilon_{1,n}>q_{2,n}/\beta_n)=p_{1,n}.
\tag{16}
\]
Thus
\[
E[H_n\mathbf1\{\beta_n\epsilon_{1,n}>q_{2,n}\}]
=p_{1,n}+O(p_{1,n}^2),
\qquad
E[H_n\mathbf1\{\epsilon_{2,n}>q_{2,n}\}]=\frac12p_{2,n},
\tag{17}
\]
the latter by independence. The intersection contributes at most \(p_{1,n}p_{2,n}\), and (15) makes the remaining two-small-jumps event \(o(p_{1,n}+p_{2,n})\). Therefore
\[
\Gamma_{21,n}(t_n)
=\frac{p_{1,n}+p_{2,n}/2}{p_{1,n}+p_{2,n}}+o(1)
=\frac12+\frac12\frac{p_{1,n}/p_{2,n}}{1+p_{1,n}/p_{2,n}}+o(1).
\tag{18}
\]
From (12),
\[
\frac{p_{1,n}}{p_{2,n}}
=\beta_n^{\alpha_{1,n}}\frac{c_{1,n}}{c_{2,n}}
q_{2,n}^{-\Delta_n}\{1+o(1)\}.
\tag{19}
\]
Moreover, scale locality and tail-constant balance imply \(\Delta_n\log c_{1,n}\to0\), and hence
\[
\Delta_n\log a_n=\frac{\lambda_n}{\alpha_{1,n}}+o(1),
\qquad
\Delta_n\log b_n=\frac{\lambda_n}{\alpha_{2,n}}+o(1).
\tag{20}
\]
Both limits equal \(\lambda/\alpha_0\). By (10) and \(\Delta_n\to0\),
\[
\Delta_n\log q_{2,n}\to\frac{\lambda}{\alpha_0}.
\tag{21}
\]
Substitution in (19), together with coefficient and tail-constant convergence, yields
\[
\frac{p_{1,n}}{p_{2,n}}\to
R(\lambda)=\beta_0^{\alpha_0}c_0e^{-\lambda/\alpha_0}.
\tag{22}
\]
Equations (18) and (22) prove the stated \(\Gamma_{21,n}\) limit.

Finally, the \((1-1/t_n)\)-quantile of \(X_{1,n}=\epsilon_{1,n}\) is \(q_{1,n}=U_{1,n}\) by continuity. On \(X_{1,n}>q_{1,n}\), positivity gives \(X_{2,n}>\beta_nU_{1,n}\), whence
\[
0\leq1-\Gamma_{12,n}(t_n)
\leq P(X_{2,n}>\beta_nU_{1,n}).
\tag{23}
\]
By (5), both \(U_{1,n}/2\) and \(\beta_nU_{1,n}/(2U_{2,n})\) lie in fixed compact subsets of \((0,\infty)\). The union bound and uniform tail normalization therefore give
\[
P(X_{2,n}>\beta_nU_{1,n})
\leq P(\epsilon_{1,n}>U_{1,n}/2)
+P(\epsilon_{2,n}>\beta_nU_{1,n}/2)
=O(t_n^{-1})\to0.
\tag{24}
\]
Thus \(\Gamma_{12,n}(t_n)\to1\). Subtracting (18) from this limit according to \(\mathrm{def{:}finite\text{-}threshold\text{-}ctc}\), and using (22), proves the asserted limit for \(D_n(t_n)\).